\section{State-Level Analysis: VA, NC, and TX}
\label{sec:state-analysis}

\subsection{Selection Rationale}

We focus on Virginia, North Carolina, and Texas because they host the three largest concentrations of active-duty military population among the 50 states:

\begin{table}[h]
\centering
\begin{tabular}{lccc}
\toprule
\textbf{State} & \textbf{Active Duty (2020)} & \textbf{Dependents} & \textbf{Congressional Districts} \\
\midrule
California & 152,000 & 113,000 & 52 \\
Virginia & 127,000 & 118,000 & 11 \\
Texas & 135,000 & 143,000 & 38 \\
North Carolina & 99,000 & 96,000 & 14 \\
Georgia & 62,000 & 52,000 & 14 \\
\midrule
National total & 1,387,000 & 1,022,000 & 435 \\
\bottomrule
\end{tabular}
\caption{Active-duty military and dependent population by state, 2020. Virginia, Texas, and North Carolina are selected for case studies due to high military concentration relative to congressional delegation size. California has more total military but proportionally less concentration. Source: Defense Manpower Data Center (DMDC) September 2020 report.}
\label{tab:state-military}
\end{table}

California has the most military personnel but also the most congressional districts (52); the dilution per district is relatively low. Virginia's 127,000 active-duty population is concentrated in Hampton Roads (Norfolk Naval Station, Langley AFB, Fort Eustis, Quantico) and creates substantial district-level inflation in a smaller delegation. North Carolina's Fort Liberty (Cumberland County) and Camp Lejeune (Onslow County) create the most extreme single-district concentration among major states.

\subsection{Virginia: Hampton Roads Concentration}

Virginia's military population is concentrated in two clusters: the Hampton Roads metro area (Norfolk Naval Station, the world's largest naval installation; Naval Station Little Creek; Langley AFB; Joint Base Langley-Eustis; Oceana NAS) and the Northern Virginia/D.C. suburbs (Quantico, Pentagon-area personnel).

The Hampton Roads cluster encompasses census tracts in Virginia Beach, Norfolk, Chesapeake, Hampton, and Newport News with military Group Quarters (type 210) populations ranging from 2,000 to 8,000 per tract. The current VA-2 congressional district (Hampton Roads area) contains approximately 48,000 active-duty and dependent residents in its Group Quarters population, plus an estimated additional 60,000 military-affiliated persons living off-base in the district. This total of approximately 108,000 military-affiliated residents represents roughly 14\% of the district's 760,000-person target population.

Without this military concentration, the district would need to expand geographically to reach the target population---absorbing more of the rural Southside Virginia counties, which would alter its partisan composition. The military population effectively compresses VA-2 geographically, concentrating it in the dense Hampton Roads core.

\subsection{North Carolina: Fort Liberty and the Cumberland County Effect}

North Carolina presents the most analytically clear case of military district inflation because the state's primary installation, Fort Liberty (formerly Fort Bragg, renamed 2023), is located in Cumberland County, which is neither a major urban center nor a rural county but a mid-sized metro ($\sim$335,000 population) whose demographic character is largely defined by the installation.

\begin{table}[h]
\centering
\begin{tabular}{lcc}
\toprule
\textbf{Population Category} & \textbf{Cumberland County (2020)} & \textbf{Share of Total} \\
\midrule
Total Census population & 335,509 & 100.0\% \\
Active-duty military (on-base GQ) & 48,752 & 14.5\% \\
Active-duty military (off-base est.) & 28,000 & 8.3\% \\
Military dependents (est.) & 51,000 & 15.2\% \\
\midrule
Military-affiliated total (est.) & 127,752 & 38.1\% \\
Civilian non-military & 207,757 & 61.9\% \\
\bottomrule
\end{tabular}
\caption{Cumberland County (NC) population breakdown. Approximately 38\% of the county's Census population is military-affiliated, making it one of the most installation-dense counties in the country.}
\label{tab:cumberland}
\end{table}

The 9th Congressional District (post-2022 map, covering Cumberland and Hoke counties plus portions of adjacent counties) is the primary Fort Liberty district. We estimate that the military population inflation in this district is approximately 6--8\% relative to the civilian-only counterfactual:

\begin{itemize}
    \item District total population target: $\sim$760,000
    \item Estimated military-affiliated district residents: 57,000--62,000 (including off-base)
    \item Military-affiliated as \% of district: 7.5--8.2\%
    \item Estimated civilian-only district population shortfall: $\sim$59,000 people
    \item Required geographic expansion if only civilians counted: $\sim$8\% more territory
\end{itemize}

The 8\% inflation is politically significant because Cumberland County borders Robeson County (majority Native American and Black; reliably Democratic) and Harnett County (majority white; reliably Republican). A district expanded to compensate for military population removal would need to absorb more of one of these counties, altering the racial and partisan composition in a measurable direction.

\subsection{Texas: Diffuse Concentration}

Texas has the largest absolute military population of any state (approximately 135,000 active-duty), but the population is spread across multiple large installations: Fort Cavazos (formerly Fort Hood, Bell County, $\sim$36,000), Joint Base San Antonio-Lackland (Bexar County, $\sim$22,000), Naval Air Station Corpus Christi (Nueces County, $\sim$4,000), Fort Bliss (El Paso County, $\sim$30,000), Dyess AFB (Taylor County, $\sim$5,000), and others.

The diffuse distribution means that no single Texas congressional district is as heavily inflated as NC-9. Fort Cavazos (Bell County) is the most concentrated, with approximately 36,000 active-duty on-base personnel in a county of roughly 370,000. The district containing Bell County (TX-31, under 2022 map) has a military-affiliated population of approximately 55,000, representing about 7\% of its congressional population target.

Texas's 38 congressional districts average the military effect broadly enough that no district shows the single-district impact seen in North Carolina. However, the combination of military inflation and the large non-citizen population (covered in N.0) means that Texas congressional maps are doubly affected by population-counting choices.

\begin{table}[h]
\centering
\begin{tabular}{llccc}
\toprule
\textbf{Installation} & \textbf{County} & \textbf{Active Duty} & \textbf{County Pop} & \textbf{Military \%} \\
\midrule
Fort Cavazos & Bell & 36,000 & 370,942 & 9.7\% \\
Ft. Bliss & El Paso & 30,000 & 865,657 & 3.5\% \\
Lackland AFB & Bexar & 22,000 & 2,009,324 & 1.1\% \\
Dyess AFB & Taylor & 5,200 & 138,034 & 3.8\% \\
\bottomrule
\end{tabular}
\caption{Major Texas military installations by active-duty population and county concentration. Fort Cavazos (Bell County) shows the highest concentration at nearly 10\% of county population.}
\label{tab:texas-installations}
\end{table}
